\labsection{3}{Static routing and the backup route}{sec:lab03-static}

\begin{labproject}{Three routers, three networks, and one deliberate failure}
\textbf{Coverage.} Chapter~\ref{ch:static-routing}.\\
\textbf{Goal.} Make a router forward to a network it is not attached to, then prove a backup route takes over when the primary fails.

\medskip\noindent\textbf{Part A --- Why the first ping fails.}\par
Configure the interfaces, then ping from R1's LAN to R3's LAN before adding any routes. It will fail, and understanding exactly why is the point of the lab.

\begin{keyidea}
A router forwards only to networks in its routing table. Directly connected networks appear automatically; everything else must be learned --- statically here, dynamically in Lab~\ref{sec:lab04-ospf}.

R1 has no entry for R3's LAN, so it drops the packet and reports \emph{Destination host unreachable}. The packet never reaches R2. Nothing is broken; the router is doing exactly what it was told.
\end{keyidea}

Confirm before configuring:

\begin{lstlisting}[style=dcnterminal]
<R1> display ip routing-table
\end{lstlisting}

Count the entries. Every one should be \cmd{Direct}. Record that output --- it is the "before" half of your evidence.

\medskip\noindent\textbf{Part B --- Static route syntax.}\par

\begin{lstlisting}[style=dcnterminal]
ip route-static <destination> <mask> <next-hop>
\end{lstlisting}

The next hop must be an address R1 can already reach --- normally the neighbouring router's interface on the shared link. A next hop R1 cannot resolve produces a route that looks correct in the table and silently discards traffic.

\begin{pitfall}
The destination is a \emph{network} address, not a host address.

\cmd{ip route-static 192.168.3.0 24 10.0.12.2} is a route to a network.\\
\cmd{ip route-static 192.168.3.1 24 10.0.12.2} is a typo that VRP will accept, normalise, and then behave unexpectedly.

Always write the network address with host bits zeroed.
\end{pitfall}

\medskip\noindent\textbf{Part C --- Routes are one-directional.}\par
Add routes on R1 so it can reach R3's LAN. Ping again. It still fails.

\begin{keyidea}
This is the single most common static-routing error, and it is worth meeting deliberately.

Your ping now reaches R3. R3 has no route back to R1's LAN, so the \emph{reply} is dropped. From R1 the symptom looks identical to the original failure, which is why students often conclude their route did not work.

Connectivity requires a route in both directions. Always configure and verify both ends before concluding anything.
\end{keyidea}

Prove it with \cmd{tracert} from R1: the trace advances further than before, then stops. That is the return path failing, not the forward path.

\medskip\noindent\textbf{Part D --- Backup routes and preference.}\par
Add a second route to the same destination through a different neighbour, with a worse preference:

\begin{lstlisting}[style=dcnterminal]
ip route-static 192.168.3.0 24 10.0.12.2              # primary, preference 60
ip route-static 192.168.3.0 24 10.0.13.3 preference 100  # backup
\end{lstlisting}

\begin{keyidea}
\term{Preference} (Huawei's term for administrative distance) ranks \emph{sources} of routing information. Lower wins. Static routes default to 60; OSPF internal routes are 10.

Only the preferred route is installed in the forwarding table. The backup sits inactive until the primary's next hop becomes unreachable, at which point it is installed automatically.

This is failover, not load sharing. Traffic does not split across the two paths --- give them equal preference if that is what you want.
\end{keyidea}

\medskip\noindent\textbf{Part E --- Prove the failover.}\par

\begin{netsteps}
  \item \cmd{display ip routing-table} --- confirm only the primary is present, with preference 60.
  \item Start a long ping from R1's LAN so you can watch the interruption.
  \item Shut down the primary link: \cmd{interface GigabitEthernet 0/0/0} then \cmd{shutdown}.
  \item \cmd{display ip routing-table} again --- the backup should now be installed.
  \item Count how many pings were lost during the transition.
  \item \cmd{undo shutdown} and confirm the primary returns.
\end{netsteps}

That lost-ping count is the real result of this lab. Static failover is not instant: the router must first decide the next hop is unreachable.

\begin{troubleshoot}
\textbf{Route shows in the table but traffic fails.} The next hop is probably unreachable. Ping the next hop itself from the same router.

\textbf{Backup never activates.} Shutting an interface on the far end may leave the local interface up, so the next hop still appears valid. Shut the local interface, or verify with \cmd{display ip routing-table protocol static}.

\textbf{Both routes appear active.} They have equal preference, so you have configured load sharing rather than failover.
\end{troubleshoot}

\begin{tryit}
Replace both static routes with one default route, \cmd{ip route-static 0.0.0.0 0 10.0.12.2}. Does connectivity still work? What have you given up, and in what topology would that choice be dangerous?
\end{tryit}
\end{labproject}

\begin{verification}
\begin{itemize}
  \item Routing tables captured before routes, after routes, and after the failover.
  \item Both directions configured, with the one-way failure demonstrated rather than skipped.
  \item Preference values visible in the table output.
  \item Packet loss during failover counted and reported.
\end{itemize}
\end{verification}

\begin{labdeliverable}
Submit configurations for all three routers, the three routing-table captures, evidence of the one-way failure and its fix, and the measured packet loss during failover with one sentence explaining why it is not zero.
\end{labdeliverable}
